\documentclass[10pt,a4paper]{amsart}
\usepackage[margin=19mm]{geometry}
\usepackage{amsmath,amssymb,mathtools}
\usepackage[scr=rsfs,cal=euler]{mathalfa}
\usepackage{xfrac}
\usepackage[unicode,hidelinks,bookmarks=false]{hyperref}
\newcommand{\ie}{i.e.\ }
\newcommand{\cf}{cf.\ }
\newcommand{\resp}{respectively\ }
\newcommand{\Subsection}{\subsection}
\providecommand{\nobreakdash}{\nobreak\hbox{-}}
\setlength{\emergencystretch}{2em}
\multlinegap0pt
\usepackage{amssymb}
\usepackage{tcolorbox}
\usepackage{dsfont}
\usepackage{xcolor}
\newtheorem{definition}{Definition}
\newtheorem{rem}{Remark}
\newtheorem{lemma}{Lemma}
\renewcommand{\P}{{\bf P}}
\newcommand{\RR}{\mathbb{R}}
\renewcommand{\d}{\, \mathrm{d} }
\newcommand{\der}{\mathrm{d}}
\DeclareMathOperator{\dv}{div}
\newcommand{\e}{{\bf e}}
\DeclareMathOperator{\essinf}{ess \, inf}
\newcommand{\p}{{\bf p}}
\newcommand{\per}{\mathrm{per}}
\DeclareMathOperator{\supp}{\mathrm{supp}}
\title{Regularity and Uniqueness for a Model of Active Particles with Angle-Averaged Diffusions}
\author{Luca Alasio, Simon Schulz}
\date{Source: arXiv:2501.11488v3 (CC BY 4.0)}
\begin{document}
\maketitle

\noindent\emph{Source and licence.} Regularity and Uniqueness for a Model of Active Particles with Angle-Averaged Diffusions. Version arXiv:2501.11488v3 (CC BY 4.0), distributed by arXiv under the Creative Commons Attribution 4.0 International licence, http://creativecommons.org/licenses/by/4.0/, as recorded in the arXiv licence metadata for this exact version and shown on the abstract page https://arxiv.org/abs/2501.11488v3. Full licence notice: \texttt{licenses/arxiv-2501.11488-CC-BY-4.0.txt}.\par\smallskip

\noindent\textit{Editorial scope.} Counted unit: the article's lemma p (source0.tex lines 510--512). The article's complete original proof of it is delivered with it (source0.tex lines 513--625, one \texttt{proof} environment of 113 lines). No mathematics is written, repaired or re-derived: the statement and the proof are the article's own lines, in the article's own notation, and no retained line is substituted or re-flowed.  Retained uncounted as the source-local context the proof needs: lines 212--215; lines 217--219; lines 254--256; lines 291--299; lines 301--303; lines 490--492; lines 503--505.  Nothing was omitted from any retained span, so the package carries no omission record.  No retained line contains a citation, so the wrapper carries no bibliography and no bibliography entry.  No retained line contains a figure environment, an image-inclusion command or drawing code, so no image is delivered and the package declares no asset dependency.  Editorial formatting only, in addition: an AMS \texttt{amsart} preamble that restates the article's own declarations and macros used by the retained lines, namely the environments \texttt{definition}, \texttt{rem}, \texttt{lemma} on separate counters, together with the standard packages those lines need.  The retained environments print in this document as Definition 1, Remark 1, Lemma 1 in order of appearance, whereas the article prints the same items with the numbering of its own full document.  The complete native source of the article as distributed in the arXiv e-print for this exact version is bundled unchanged as \texttt{sources/171853/source0.tex}.
\begin{equation}\label{eq:main eqn}
    \partial_t f +   \dv \!\big(  (1- \rho) f \e(\theta)\big)
= \dv \!\big( (1-\rho)\nabla f + f \nabla \rho \big) + \partial_{\theta}^2 f, 
\end{equation}

\begin{equation}\label{eq:rho def}
    \rho(t,x) := \int_0^{2\pi} \!\!\! f(t,x,\theta) \d \theta.
\end{equation}

\begin{equation}\label{eq:polarisation def}
    \p(t,x) := \int_0^{2\pi} \!\!\! f(t,x,\theta) \, \e(\theta) \d \theta. 
\end{equation}

\begin{definition}[Function Space $\mathcal{X}$]\label{def:function space}
   We define the function space $\mathcal{X}$ to be the family of functions $f \in L^3(0,T;L^3_\per(\Upsilon))$ which satisfy, with $\rho := \int_0^{2\pi} f \d \theta$, the following: 
    \begin{itemize}
        \item[(i)] $f \geq 0$ a.e.~in $\Upsilon_T$; 
        \item[(ii)] $\partial_\theta \sqrt{f} \in L^2(\Upsilon_T)$ and $\partial_t f \in (L^6(0,T;W^{1,6}_\per(\Upsilon))'$; 
        \item[(iii)] $0 \leq \rho \leq 1$ a.e.~in $\Omega_T$, $\partial_t \rho \in (L^2(0,T;H^1_\per(\Omega))'$, and $\nabla\sqrt{1-\rho} \in L^2(\Omega_T)$; 
        \item[(iv)] $\sqrt{1-\rho}\nabla\sqrt{f} \in L^2(\Upsilon_T)$. 
    \end{itemize}
\end{definition}

\begin{equation}\label{eq:early boundedness nabla f weighted in intro}
    \sqrt{1-\rho}\nabla f , \partial_\theta f \in L^{\frac{3}{2}}(\Upsilon_T). 
\end{equation}

\begin{equation}\label{eq:polarisation eq}
    \partial_t \p + \dv((1-\rho)\P) = \dv((1-\rho)\nabla \p + \p \otimes \nabla \rho) - \p 
\end{equation}

\begin{rem}[Weak formulation of polarisation equation]\label{rem:sense of p eqn to begin with}
    All terms in \eqref{eq:polarisation eq} are well-defined since $(1-\rho)\nabla \p \in L^{\frac{3}{2}}(\Omega_T)$ from \eqref{eq:early boundedness nabla f weighted in intro} and $\p \otimes \nabla \rho \in L^2(\Omega_T)$ from Definition \ref{def:function space}. By the previous bounds, \eqref{eq:polarisation eq} holds in duality to $L^3(0,T;(W^{1,3}_\per(\Omega))')$, $\partial_t \p \in L^{\frac{3}{2}}(0,T;(W^{1,3}_\per(\Omega))')$. 
\end{rem}

\begin{lemma}[$H^1$-type estimate for $\p$]\label{lem:H1 p}
  Let $f$ be a weak solution of \eqref{eq:main eqn}, with $\rho$ and $\p$ as per \eqref{eq:rho def}--\eqref{eq:polarisation def}. Then,  if for some $t\geq 0$ it holds $\essinf_{\Omega_{t,T}}(1-\rho)>0$, then there holds  $\nabla \p \in L^2(\Omega_{t,T})$, $\partial_t \p \in L^2(t,T;(H^1_\per(\Omega))')$. 
\end{lemma}

\begin{proof}
    The underlying idea is to test the equation \eqref{eq:polarisation eq} with $\p$ itself. Formally, 
    \begin{equation*}
    \begin{aligned}
         \frac{1}{2}\frac{\der}{\der t}\!\int_\Omega\! |\p|^2 \d x \!+\! \frac{1}{2}\int_\Omega\! (1\!-\!\rho) |\nabla \p|^2 \d x \!\leq\! |\Omega| \!+\! \Vert \P \Vert_{L^\infty(\Omega_T)} \! \int_\Omega\! |\nabla \sqrt{1-\rho}|^2 \d x \!-\! \frac{1}{2}\int_\Omega |\p|^2 \d x, 
    \end{aligned}
\end{equation*}
and integrating in time yields the sought estimate. However, this argument is not rigorous since \emph{a priori} we do not know that $\p$ is an admissible test function. To overcome this, we test the equation with a smoothed version of $\p$ and conclude with a limiting argument. 

    \smallskip 

    \noindent 1. \textit{Mollification}: Let $\eta$ be the usual Friedrichs bump function, and $\{\eta^\varepsilon\}_{\varepsilon>0}$ the usual sequence of Friedrichs mollifiers in $\mathbb{R}^2$; we only mollify with respect to $x$. In what follows, we use the notations $\p^\varepsilon = \eta^\varepsilon*\p$, $D = (1-\rho)\P$, and $D^\varepsilon = \eta^\varepsilon*D$. Since $\eta^\varepsilon$ is an admissible test function to insert into the weak formulation of \eqref{eq:polarisation eq}, we mollify the equation: 
    \begin{equation}\label{eq:mollified p eqn}
        \partial_t \p^\varepsilon + \dv D^\varepsilon = \dv((1-\rho)\nabla\p^\varepsilon + \p^\varepsilon \otimes \nabla \rho) - \p^\varepsilon + \dv (E^\varepsilon + F^\varepsilon), 
    \end{equation}
    where the error terms $E^\varepsilon$ and $F^\varepsilon$ are given by 
    \begin{equation}\label{eq:error terms mollified p eqn}
        E^\varepsilon = \eta^\varepsilon*((1-\rho)\nabla \p) - (1-\rho)\eta^\varepsilon*\nabla \p , \qquad F^\varepsilon = \eta^\varepsilon*(\p \otimes \nabla \rho) - (\eta^\varepsilon*\p )\otimes \nabla \rho. 
    \end{equation}

    \smallskip 

    \noindent 2. \textit{Estimates on drift term and error terms}: We estimate the error terms $D^\varepsilon, E^\varepsilon, F^\varepsilon$. Our aim is to show the following: there exists a positive constant $C$ independent of $\varepsilon$ such that 
    \begin{equation}\label{eq:unif bounds for H1 for p with mollifications}
        \Vert D^\varepsilon \Vert_{L^\infty(\Omega_T)} + \Vert E^\varepsilon \Vert_{L^2(\Omega_T)} + \Vert F^\varepsilon \Vert_{L^2(\Omega_T)} \leq C. 
    \end{equation}
    Once the above is established, we will be able to proceed as per the formal argument presented at the start of the proof. First, $\Vert D^\varepsilon (t,\cdot) \Vert_{L^\infty(\Omega)} \leq C \Vert D(t,\cdot) \Vert_{L^\infty(\Omega)} \leq C \Vert \P \Vert_{L^\infty(\Omega_T)} \leq C,$ for some positive constant $C$ independent of $\varepsilon$, where we used standard properties of the mollifier to establish the first inequality. It follows from taking the essential supremum in the time variable that, for some positive constant independent of $\varepsilon$, 
    \begin{equation}\label{eq:first error term H1 proof p}
        \Vert D^\varepsilon \Vert_{L^\infty(\Omega_T)} \leq C. 
    \end{equation}
Next, the error term $F^\varepsilon$ is controlled using elementary properties of the mollifier, as follows: 
\begin{equation*}
    \Vert \eta^\varepsilon*(\p \otimes \nabla \rho)(t,\cdot) \Vert_{L^2(\Omega)} \leq C \Vert (\p \otimes \nabla \rho) (t,\cdot) \Vert_{L^2(\Omega)} \leq C \Vert \p \Vert_{L^\infty(\Omega_T)} \Vert \nabla \rho(t,\cdot) \Vert_{L^2(\Omega)}, 
\end{equation*}
for a.e.~$t \in(0,T)$, for some positive constant $C$ independent of $\varepsilon$, and similarly 
\begin{equation*}
    \Vert (\eta^\varepsilon*\p )\otimes \nabla \rho (t,\cdot) \Vert_{L^2(\Omega)} \!\leq\! \Vert (\eta^\varepsilon*\p )(t,\cdot) \Vert_{L^\infty(\Omega)}\Vert \nabla \rho(t,\cdot) \Vert_{L^2(\Omega)} \!\leq\! C \Vert \p \Vert_{L^\infty(\Omega_T)} \Vert \nabla \rho(t,\cdot) \Vert_{L^2(\Omega)}, 
\end{equation*}
whence, by integrating in time, there exists a positive $C$ independent of $\varepsilon$ such that 
\begin{equation}\label{eq:second error term H1 proof p}
    \Vert F^\varepsilon \Vert_{L^2(\Omega_T)} \leq C. 
\end{equation}
For $E^\varepsilon$, write $E^\varepsilon(t,x) = \int_{\mathbb{R}^2} \eta^\varepsilon(x-y)\big[ \rho(t,x)-\rho(t,y) \big]\nabla\p(t,y) \d y$. Integrating by parts, 
\begin{equation*}
   \begin{aligned}
       E^\varepsilon(t,x) =\!\!  \underbrace{\int_{\mathbb{R}^2} \!\!\p(t,y) \otimes \nabla \eta^\varepsilon(x-y) \big[ \rho(t,x)-\rho(t,y) \big] \d y}_{=: E^\varepsilon_1(t,x)} + \!\underbrace{\int_{\mathbb{R}^2} \!\!\eta^\varepsilon(x-y) \p(t,y) \otimes \nabla \rho(t,y) \d y}_{=: E^\varepsilon_2(t,x)} . 
   \end{aligned} 
\end{equation*}
Using standard properties of mollifiers, $\Vert E^\varepsilon_2(t,\cdot) \Vert_{L^2(\Omega)} \leq C \Vert \nabla \rho(t,\cdot) \Vert_{L^2(\Omega)} \Vert \p \Vert_{L^\infty(\Omega_T)},$ hence 
\begin{equation}\label{eq:Eeps2 bound unif}
    \Vert E^\varepsilon_2 \Vert_{L^2(\Omega_T)} \leq C \Vert \nabla \rho \Vert_{L^2(\Omega_T)} \Vert \p \Vert_{L^\infty(\Omega_T)}, 
\end{equation}
for some positive constant $C$ independent of $\varepsilon$. Meanwhile, note that, for a.e.~$x,y\in \RR^2$,
\begin{equation*}
   \begin{aligned} 
  \rho(t,x)-\rho(t,y) = \int_0^1 \frac{\der}{\der \sigma}\rho(t,\sigma x +(1-\sigma)y)  \d \sigma = \int_0^1 \nabla \rho(t,\sigma x +(1-\sigma)y) \!\cdot\! (x-y)  \d \sigma, 
   \end{aligned}
\end{equation*}
and since $\supp \eta^\varepsilon = B(0,\varepsilon)$, it follows that, for all $y \in \supp \eta^\varepsilon(x-\cdot)$, there holds 
\begin{equation}\label{eq:rho diff for mollification argument}
   \begin{aligned} 
  |\rho(t,x)-\rho(t,y)| \leq \varepsilon \int_0^1 |\nabla \rho(t,\sigma x +(1-\sigma)y)| \d \sigma. 
   \end{aligned}
\end{equation}
In turn, using the uniform boundedness of $\p$, there holds 
\begin{equation*}
    \begin{aligned}
        E^\varepsilon_1(t,x) &\leq \varepsilon^{-2} \int_{B(x,\varepsilon)} \big|\nabla \eta(\frac{x-y}{\varepsilon}) \big| \bigg(  \int_0^1 |\nabla \rho(t,\sigma x +(1-\sigma)y)| \d \sigma\bigg) \d y \\ 
        &= \int_{B(0,1)} |\nabla \eta (w)| \bigg(  \int_0^1 |\nabla \rho \big(t, x - (1-\sigma)\varepsilon w \big)| \d \sigma\bigg) \d w, 
    \end{aligned}
\end{equation*}
and, by Jensen's inequality, we get $|E^\varepsilon_1(t,x)|^2 \leq C \int_{B(0,1)} \int_0^1 |\nabla \rho \big(t, x - (1-\sigma)\varepsilon w \big)|^2 \d \sigma \d w$, for some positive constant $C$ independent of $\varepsilon$. An application of the Fubini--Tonelli Theorem then implies, using periodicity of $\nabla \rho$ and translation invariance of Lebesgue measure, 
\begin{equation*}
    \Vert E^\varepsilon_1(t,\cdot) \Vert^2_{L^2(\Omega)} \leq C \int_{B(0,1)} \int_0^1 \int_\Omega |\nabla \rho(t,x)|^2 \d x \d \sigma \d w \leq C \Vert \nabla \rho(t,\cdot) \Vert^2_{L^2(\Omega)}. 
\end{equation*}
By integrating in time, and combining with \eqref{eq:Eeps2 bound unif}, we deduce $\Vert E^\varepsilon \Vert_{L^2(\Omega_T)} \leq C$, independently of $\varepsilon$. Combining this bound with \eqref{eq:first error term H1 proof p} and \eqref{eq:second error term H1 proof p}, we obtain \eqref{eq:unif bounds for H1 for p with mollifications}. 

\smallskip 

    \noindent 3. \textit{$H^1$-type estimate}: Note that $\p^\varepsilon \in L^\infty(0,T;C^\infty_{\per}(\Omega))$, whence, in view of Remark \ref{rem:sense of p eqn to begin with}, $\p^\varepsilon$ is an admissible test function to insert in the weak formulation of \eqref{eq:polarisation eq}. Testing the equation with $\p^\varepsilon$, using the uniform bounds \eqref{eq:unif bounds for H1 for p with mollifications} and Young's inequality, we get 
    \begin{equation*}
        \begin{aligned}
            \frac{1}{2}\frac{\der}{\der t}\int_\Omega |\p^\varepsilon|^2 \d x + \essinf_{\Omega_{t,T}}(1-\rho) \int_\Omega |\nabla \p^\varepsilon|^2 \d x \leq C_{\essinf_{\Omega_{t,T}}(1-\rho)} + \Vert \nabla \sqrt{1-\rho(t,\cdot)} \Vert^2_{L^2(\Omega)} , 
        \end{aligned}
    \end{equation*}
with $C$ independent of $\varepsilon$. Integrating in time, as $\Vert \p^\varepsilon(0,\cdot)\Vert_{L^\infty(\Omega)} \leq C \Vert \p(0,\cdot)\Vert_{L^\infty(\Omega)} \leq C$, 
    \begin{equation}\label{eq:just before H1 final bound for p}
        \begin{aligned}
            \Vert \p^\varepsilon\Vert_{L^\infty(t,T;L^2(\Omega))}^2 + \Vert \nabla \p^\varepsilon\Vert_{L^2(\Omega_{t,T})}^2 &\leq C\big( 1+ \Vert \nabla \sqrt{1-\rho} \Vert^2_{L^2(\Omega_T)}). 
        \end{aligned}
    \end{equation}
    We take the limit $\varepsilon \to 0$, and the result follows from the lower semicontinuity of the norms. 
    
    \iffalse The product rule implies $\sqrt{1-\rho}\nabla \p^\varepsilon = \nabla(\sqrt{1-\rho}\p^\varepsilon) - \p^\varepsilon \otimes \nabla\sqrt{1-\rho}$, from which we deduce, by \eqref{eq:just before H1 final bound for p}, that $\{\sqrt{1-\rho}\p^\varepsilon \}_\varepsilon$ is uniformly bounded in $L^2(0,T;H^1(\Omega))$. It follows from Alaoglu's Theorem and the strong convergence $\p^\varepsilon \to \p$ a.e.~$\Omega_T$ and in $L^q(\Omega_T)$ for all $q \in [1,\infty)$ that there holds 
    \begin{equation}\label{eq:sidestep to take limit of sqrt p}
        \sqrt{1-\rho}\p^\varepsilon \rightharpoonup \sqrt{1-\rho}\p \quad \text{weakly in } L^2(0,T;H^1(\Omega)). 
    \end{equation}
    Hence, using \eqref{eq:just before H1 final bound for p}, the triangle inequality, and the weak lower semicontinuity of the norm, 
    \begin{equation*}
       \begin{aligned}
         \Vert \sqrt{1-\rho}\nabla \p \Vert_{L^2(\Omega_T)} &\leq \Vert \p \otimes \nabla\sqrt{1-\rho} \Vert_{L^2(\Omega_T)} + \Vert \nabla(\sqrt{1-\rho}\p) \Vert_{L^2(\Omega_T)} \\ 
         &\leq \Vert \nabla\sqrt{1-\rho} \Vert_{L^2(\Omega_T)} + \liminf_{\varepsilon\to 0} \Vert \nabla(\sqrt{1-\rho}\p^\varepsilon) \Vert_{L^2(\Omega_T)} \\ 
         &\leq C\big( 1+ \Vert \nabla \sqrt{1-\rho} \Vert_{L^2(\Omega_T)} \big), 
       \end{aligned}
    \end{equation*}
and therefore $\sqrt{1-\rho}\nabla \p \in L^2(\Omega_T)$. \fi 



    \smallskip 

    \noindent 4. \textit{Boundedness of time derivative}: We obtain $\partial_t \p \in L^2(t,T;(H^1(\Omega))')$ directly from the equation, using the boundedness $\sqrt{1-\rho}\nabla \p \in L^2(\Omega_T)$ and $\nabla \rho \in L^2(\Omega_T)$. 
\end{proof}

\end{document}
